\documentclass{article}

\usepackage{amsmath,amssymb}
\usepackage{xurl}
\usepackage[hidelinks]{hyperref}
\setlength{\emergencystretch}{2em}

% Shared commands for notes in docs/paper-gaps/.

\DeclareUnicodeCharacter{2080}{\ifmmode{}_{0}\else\textsubscript{0}\fi}
\DeclareUnicodeCharacter{2081}{\ifmmode{}_{1}\else\textsubscript{1}\fi}
\DeclareUnicodeCharacter{2082}{\ifmmode{}_{2}\else\textsubscript{2}\fi}
\DeclareUnicodeCharacter{2099}{\ifmmode{}_{n}\else\textsubscript{n}\fi}

\newcommand{\Lean}{\textsc{Lean}}
\ExplSyntaxOn
\NewDocumentCommand{\leanid}{m}{
  \ifmmode
    \textsf{\detokenize{#1}}
  \else
    \group_begin:
    \urlstyle{sf}
    \tl_set:Nn \l_tmpa_tl {#1}
    \tl_replace_all:Nnn \l_tmpa_tl {\_} {_}
    \exp_args:NV \path \l_tmpa_tl
    \group_end:
  \fi
}
\ExplSyntaxOff
\providecommand{\tr}{\operatorname{tr}}
\newcommand{\tnleanIssueBase}{https://github.com/LionSR/TNLean/issues/}
\newcommand{\tnleanPrBase}{https://github.com/LionSR/TNLean/pull/}
\newcommand{\tnleanDocBase}{https://sirui-lu.com/TNLean/docs/}
\newcommand{\ghissue}[1]{\href{\tnleanIssueBase#1}{GitHub issue \##1}}
\newcommand{\ghpr}[1]{\href{\tnleanPrBase#1}{PR \##1}}
\newcommand{\doclink}[1]{\href{\tnleanDocBase#1.html}{\path{#1}}}

% Dirac notation and the identity operator, as in blueprint/src/macros/common.tex.
% \providecommand keeps note-local definitions authoritative.
\usepackage{dsfont}
\providecommand{\ket}[1]{|#1\rangle}
\providecommand{\bra}[1]{\langle #1|}
\providecommand{\braket}[2]{\langle #1 | #2 \rangle}
\providecommand{\ketbra}[2]{|#1\rangle\!\langle #2|}
\providecommand{\Id}{\mathds{1}}
\providecommand{\one}{\mathds{1}}
\providecommand{\C}{\mathbb{C}}
\providecommand{\MN}[1]{M_{#1}(\C)}
\providecommand{\MD}{M_D(\C)}

% External referencing for paper sources.  The build helper writes sanitized
% auxiliary files under build/paper-aux/ and excludes duplicated source labels.
\usepackage{xr}

% Import a generated paper auxiliary file with a caller-supplied label prefix.
% The candidate paths support builds from docs/paper-gaps/, the repository root,
% and build/paper-aux/ itself.
\newcommand{\importpaperaux}[2]{%
  \IfFileExists{../../build/paper-aux/#2.aux}{%
    \externaldocument[#1]{../../build/paper-aux/#2}%
  }{%
    \IfFileExists{build/paper-aux/#2.aux}{%
      \externaldocument[#1]{build/paper-aux/#2}%
    }{%
      \IfFileExists{#2.aux}{%
        \externaldocument[#1]{#2}%
      }{}%
    }%
  }%
}

% General external reference macros
\newcommand{\paperref}[2]{\ref{#1-#2}}
\newcommand{\papereqref}[2]{(\ref{#1-#2})}

% CPSV16 source (1606.00608 / MPDO-22-12-17-2.tex) external document and shortcuts
\importpaperaux{cpsv16-}{1606.00608/MPDO-22-12-17-2-tnlean}
\newcommand{\cpsvref}[1]{\paperref{cpsv16}{#1}}
\newcommand{\cpsveqref}[1]{\papereqref{cpsv16}{#1}}

% Verdict marker: \gapnote{<kind>}{<status>}. Kind is the policy
% classification (clarification, local-correction, scope-restriction,
% unfaithful, false-source, open-gap); status is open, wip (elimination
% underway), resolved, or historical. Machine-read by the site index;
% prints a verdict line.
\newcommand{\gapnote}[2]{\par\noindent\textbf{Verdict:} #1 (#2).\par}


\title{CPGSV17 Blocked-Basis Chi Families and Uniform BNT Labels}
\author{}
\date{2026-07-09}

\begin{document}
\maketitle
\gapnote{scope-restriction}{resolved}

\section{Source statement}

Theorem~4.14(ii) of
Cirac--P\'erez-Garc\'ia--Schuch--Verstraete states that there are diagonal
matrices
\begin{align}
    \chi_{\alpha,\beta,\gamma}
        &> 0
\label{eq:cpgsv17_blocked_chi_uniformity_source_positive_chi}
\end{align}
whose entries are positive and for which, at every positive length $L$, the
operators $O_L(M_\alpha)$ linearly span an algebra with structure
coefficients
\begin{align}
    c^{(L)}_{\alpha,\beta,\gamma}
        &= \tr\!\left(\chi_{\alpha,\beta,\gamma}^{L}\right).
\label{eq:cpgsv17_blocked_chi_uniformity_source_trace_power}
\end{align}
Here $M_\alpha$ is the normal tensor with BNT label $\alpha$, and
$O_L(M_\alpha)$ is its length-$L$ closed operator. The labels
$\alpha,\beta,\gamma$ are fixed BNT labels rather than basis labels chosen
separately at each length. Thus each matrix
$\chi_{\alpha,\beta,\gamma}$ in
(\ref{eq:cpgsv17_blocked_chi_uniformity_source_trace_power}) is independent
of $L$; only its exponent varies. This is the assertion of
Theorem~4.14(ii) in Ref.~\cite{Cirac2016MPDO_arXiv}.

The formal development instead chooses, for each blocked length $n$, a support
algebra $\mathcal A_n$. The basis of $\mathcal A_n$ may therefore depend on
$n$.

\section{Formalized blocked-basis analogue}

The construction in
\path{TNLean/MPS/MPDO/AlgebraStructure.lean} chooses bases of the support
algebras $\mathcal A_n$. Its multiplication coefficients have indices
\begin{align}
    i,j &\in \operatorname{Basis}(\mathcal A_n),
\label{eq:cpgsv17_blocked_chi_uniformity_blocked_source_indices}\\
    k &\in \operatorname{Basis}(\mathcal A_{2n}),
\label{eq:cpgsv17_blocked_chi_uniformity_blocked_target_index}
\end{align}
because the represented product is
\begin{align}
    \mathcal A_n\times\mathcal A_n
        &\longrightarrow \mathcal A_{2n}.
\label{eq:cpgsv17_blocked_chi_uniformity_blocked_product}
\end{align}
The declarations
\leanid{MPOTensor.AlgebraStructureData.BlockedStructureChiFamily} and
\leanid{MPOTensor.AlgebraStructureData.HasBlockedStructureChiTracePowerForm}
express the positive-length, length-dependent formula
\begin{align}
    c^{(n)}_{i,j,k}
        &= \sum_r \chi_{n,i,j,k,r}^{\,n},
\label{eq:cpgsv17_blocked_chi_uniformity_blocked_power_sum}\\
        &= \tr\!\left(\chi_{n,i,j,k}^{\,n}\right).
\label{eq:cpgsv17_blocked_chi_uniformity_blocked_trace_power}
\end{align}
In
(\ref{eq:cpgsv17_blocked_chi_uniformity_blocked_power_sum})--%
(\ref{eq:cpgsv17_blocked_chi_uniformity_blocked_trace_power}), $n$ is the
source blocking length of the two factors in $\mathcal A_n$, whereas $k$
indexes a basis of $\mathcal A_{2n}$. Both the family and its trace-power
predicate are indexed only by $n>0$. They therefore contain neither a
length-zero component nor the unintended constraint
\begin{align}
    c^{(0)}_{i,j,k}
        &= \dim\!\left(\chi_{0,i,j,k}\right).
\label{eq:cpgsv17_blocked_chi_uniformity_excluded_zero_length}
\end{align}

\section{Gap}

The blocked-basis construction is a scope restriction, not the uniform
BNT-label statement of Theorem~4.14(ii) in
Ref.~\cite{Cirac2016MPDO_arXiv}. It permits the diagonal matrix to depend on
$n$, since the available basis indices depend on $\mathcal A_n$ and
$\mathcal A_{2n}$. By contrast, the source theorem assigns one matrix to each
fixed triple of BNT labels and uses it at every positive length.

The two constructions also represent different products. The source
coefficients belong to the same-length operator algebra:
\begin{align}
    O_L(M_\alpha)O_L(M_\beta)
        &= \sum_\gamma
            c^{(L)}_{\alpha,\beta,\gamma}O_L(M_\gamma).
\label{eq:cpgsv17_blocked_chi_uniformity_source_same_length_product}
\end{align}
The formal blocked product is instead
(\ref{eq:cpgsv17_blocked_chi_uniformity_blocked_product}). Its inputs have
source length $n$, but its output is expanded in a basis at length $2n$.
Consequently, the exponent in
\begin{align}
    c^{(n)}_{i,j,k}
        &= \tr\!\left(\chi_{n,i,j,k}^{\,n}\right)
\label{eq:cpgsv17_blocked_chi_uniformity_source_length_exponent}
\end{align}
is the source length of the factors, not the codomain length. This convention
does not replace the same-length multiplication law
(\ref{eq:cpgsv17_blocked_chi_uniformity_source_same_length_product}).

Theorem~4.14(ii) also contains the length-one idempotent condition
\begin{align}
    m_\gamma
        &= \sum_{\alpha,\beta}
            c^{(1)}_{\alpha,\beta,\gamma}m_\alpha m_\beta.
\label{eq:cpgsv17_blocked_chi_uniformity_source_idempotent_condition}
\end{align}
This condition appears at lines~981--985 of
Ref.~\cite{Cirac2016MPDO_arXiv}; it is independent of the blocked-basis
trace-power predicate considered here.

It follows that
\leanid{MPOTensor.AlgebraStructureData.HasBlockedStructureChiTracePowerForm}
must not be cited as a formalization of Theorem~4.14(ii). It is a local
blocked-basis coefficient statement and becomes relevant only after the
uniform BNT-label theorem has been separated from the chosen blocked bases.

\section{Explicit vertical closure for the rescaling-stable example R}

The example-specific vertical calculation closes the one-label case for $R$.
It does not remove the general blocked-basis scope restriction or the
correction required by the printed form of Theorem~4.14.

Let $E_{ij}$ be the $2\times2$ matrix unit. Define
\begin{align}
    A^0 &= \frac45E_{00},
&
    A^1 &= \frac45E_{11},
\label{eq:cpgsv17_blocked_chi_uniformity_example_diagonal_letters}\\
    A^2 &= \frac35E_{01},
&
    A^3 &= \frac35E_{10}.
\label{eq:cpgsv17_blocked_chi_uniformity_example_offdiagonal_letters}
\end{align}
The horizontal MPO tensor $R$ has entries
\begin{align}
    R^{pq}_{ab}
        &= \frac{25}{32}
            \left(A^a\otimes\overline{A^b}\right)_{pq}.
\label{eq:cpgsv17_blocked_chi_uniformity_horizontal_tensor_R}
\end{align}
Its normalized vertical component is the $16$-letter tensor
\begin{align}
    M_0^{(a,b)}
        &= A^a\otimes\overline{A^b},
\label{eq:cpgsv17_blocked_chi_uniformity_vertical_component_M0}
\end{align}
and $O_L(M_0)$ denotes the corresponding length-$L$ closed operator. Set
\begin{align}
    W
        &= \begin{pmatrix}
            16/25 & 9/25\\
            9/25 & 16/25
        \end{pmatrix},
&
    W_L
        &= W^{\otimes L}.
\label{eq:cpgsv17_blocked_chi_uniformity_example_W_definitions}
\end{align}
The notation $W^L$ denotes an ordinary matrix power, whereas $W_L$ denotes
the $L$-fold Kronecker power.

The module
\path{TNLean/MPS/MPDO/RescalingStableExplicitVerticalBNT.lean} defines the
one-label clause
\leanid{MPOTensor.RescalingStableLengthDependentRFP.R_oneLabelBNTAlgebraTensorClause}.
Its normalized vertical component has $16$ physical letters, bond dimension
$4$, and satisfies
(\ref{eq:cpgsv17_blocked_chi_uniformity_vertical_component_M0}). The
multiplicity matrix and the distinct chi matrix are
\begin{align}
    \mu_0
        &= \begin{pmatrix}25/32\end{pmatrix},
\label{eq:cpgsv17_blocked_chi_uniformity_example_multiplicity_matrix}\\
    \chi_{0,0,0}
        &= \operatorname{diag}\!\left(1,\frac7{25}\right).
\label{eq:cpgsv17_blocked_chi_uniformity_example_chi_matrix}
\end{align}
The vertical reconstruction is
\begin{align}
    \operatorname{vertical}(R)
        &= \frac{25}{32}M_0.
\label{eq:cpgsv17_blocked_chi_uniformity_vertical_reconstruction}
\end{align}
For every positive length $L$, the component operator obeys
\begin{align}
    O_L(M_0)^2
        &= \left(1+\left(\frac7{25}\right)^L\right)O_L(M_0),
\label{eq:cpgsv17_blocked_chi_uniformity_example_operator_square}\\
        &= \tr\!\left(\chi_{0,0,0}^{L}\right)O_L(M_0).
\label{eq:cpgsv17_blocked_chi_uniformity_example_operator_chi_trace}
\end{align}
At length one,
\begin{align}
    m_0
        &= \tr(\mu_0)
         = \frac{25}{32},
\label{eq:cpgsv17_blocked_chi_uniformity_example_trace_scalar}\\
    c^{(1)}_{0,0,0}
        &= \tr(\chi_{0,0,0})
         = \frac{32}{25},
\label{eq:cpgsv17_blocked_chi_uniformity_example_length_one_coefficient}
\end{align}
and the idempotent identity becomes
\begin{align}
    \frac{25}{32}
        &= \frac{32}{25}
            \left(\frac{25}{32}\right)^2.
\label{eq:cpgsv17_blocked_chi_uniformity_example_idempotent_identity}
\end{align}
Thus
\leanid{MPOTensor.RescalingStableLengthDependentRFP.oneLabelCoeffs} and
\leanid{MPOTensor.RescalingStableLengthDependentRFP.oneLabelChi} are the
actual coefficient and diagonal data of this explicit vertical clause.

The contraction orientation matters. Let
$\operatorname{rot}(\operatorname{vertical}(R))$ be the MPO obtained by
splitting the $16$-valued physical index of the vertical tensor into two
four-valued indices. Then
\begin{align}
    \operatorname{mpo}\!\left(\operatorname{rot}(\operatorname{vertical}(R)),L\right)
        &= \left(\frac{25}{32}\right)^L O_L(M_0).
\label{eq:cpgsv17_blocked_chi_uniformity_rotated_vertical_contraction}
\end{align}
This rotated contraction is not the original horizontal tensor $R$, although
both its physical and bond dimensions equal $4$. At length one,
$\operatorname{mpo}(R,1)$ has rank two, while $O_1(M_0)$ is a rank-one
pure-state operator. Hence they are not proportional, and the identity
\begin{align}
    \operatorname{mpo}(R,L)
        &= \left(\frac{25}{32}\right)^L O_L(M_0)
\label{eq:cpgsv17_blocked_chi_uniformity_false_horizontal_vertical_identity}
\end{align}
must not be used.

The horizontal countercalculation remains the regression boundary. The
ordinary power $W^L$ differs from the Kronecker power $W_L$ in the
factorization of the original horizontal operator. The false scalar product
law would imply at $L=1$ that
\begin{align}
    \frac{25}{32}W^2
        &= \frac{32}{25}W.
\label{eq:cpgsv17_blocked_chi_uniformity_false_W_scalar_law}
\end{align}
The $(0,0)$ entries of the two sides are instead
\begin{align}
    \left(\frac{25}{32}W^2\right)_{00}
        &= \frac{337}{800},
\label{eq:cpgsv17_blocked_chi_uniformity_left_W_entry}\\
    \left(\frac{32}{25}W\right)_{00}
        &= \frac{512}{625}.
\label{eq:cpgsv17_blocked_chi_uniformity_right_W_entry}
\end{align}

This closure applies only to $R$. It neither identifies general blocked
support-algebra bases with uniform BNT labels nor proves that every
blocked-basis chi family is length independent. It also does not attach
arbitrary uniform rescalings of chi matrices to tensor presentations. The
formal algebra characterization of
Theorem~4.14(i)$\Leftrightarrow$(ii) is complete, while the printed theorem
remains partial because clause~(iii) requires the documented active-support
correction.

\section{Relation to the algebra-to-fusion converse}

The relevant source line ranges serve different purposes. Lines~1830--1922
prove Proposition~4.13, the vertical canonical-form statement, and do not
extract $\chi_{\alpha,\beta,\gamma}$. The forward construction
$(\mathrm{i})\Rightarrow(\mathrm{ii})$ constructs chi matrices at
lines~2020--2042. The converse
$(\mathrm{ii})\Rightarrow(\mathrm{i})$ begins at line~2046 and assumes both
the positive chi-coefficient formula and the idempotent identity
\cite{Cirac2016MPDO_arXiv}.

Let $E=E_1$ be the one-site transfer map. The present support-algebra
predicate gives
\begin{align}
    \operatorname{Fix}\!\left((E^n)^\dagger\right)
        &= \operatorname{Fix}(E^\dagger),
&
    n &> 0.
\label{eq:cpgsv17_blocked_chi_uniformity_fixed_point_tower}
\end{align}
This identity excludes root-of-unity peripheral eigenvalues of $E^\dagger$,
but it does not imply
\begin{align}
    E^2 &= E.
\label{eq:cpgsv17_blocked_chi_uniformity_transfer_idempotence}
\end{align}

The failure is formalized by
\leanid{MPOTensor.exists_hasBlockedAdjointFixedPointAlgebraTower_not_isZCL}.
The counterexample is the phase-flip channel with Kraus operators
\begin{align}
    K_0 &= \frac35 I,
&
    K_1 &= \frac45 Z.
\label{eq:cpgsv17_blocked_chi_uniformity_phase_flip_kraus}
\end{align}
This channel is trace preserving, fixes the positive-definite matrix $I$, and
has off-diagonal eigenvalue
\begin{align}
    \lambda_{\mathrm{off}}
        &= -\frac7{25}.
\label{eq:cpgsv17_blocked_chi_uniformity_phase_flip_eigenvalue}
\end{align}
Every positive power therefore has the same adjoint fixed-point algebra, even
though the transfer map is not idempotent. The implication from the present
support-algebra predicate is false even after adding trace preservation and a
faithful fixed point.

The converse proof of Theorem~4.14 in
Ref.~\cite{Cirac2016MPDO_arXiv} uses a stronger coefficient comparison. In
Appendix~C.4, lines~2046--2085 compare the BNT expansions from the one-site
and two-site decompositions. The source identity
\texttt{eq:algebra-appendix}, stated at lines~1933--1936 and used at
lines~2048--2050, is
\begin{align}
    \sum_\gamma \tr(\nu_\gamma^L)O_L(A_\gamma)
        &=
        \sum_{\alpha,\beta,\gamma}
        \tr\!\left((\mu_\alpha\otimes\mu_\beta
            \otimes\chi_{\alpha,\beta,\gamma})^L\right)O_L(M_\gamma).
\label{eq:cpgsv17_blocked_chi_uniformity_source_bnt_comparison}
\end{align}
Because both $A_\gamma$ and $M_\gamma$ are BNT tensors, the comparison first
gives, after relabelling,
\begin{align}
    A_\gamma
        &= e^{i\theta_\gamma}
            Z_\gamma M_\gamma Z_\gamma^{-1}.
\label{eq:cpgsv17_blocked_chi_uniformity_phase_gauge_relation}
\end{align}
Proposition~4.13 then removes the phase and identifies $Z_\gamma$ with a
unitary $U_\gamma$:
\begin{align}
    A_\gamma
        &= U_\gamma M_\gamma U_\gamma^{-1}.
\label{eq:cpgsv17_blocked_chi_uniformity_unitary_gauge_relation}
\end{align}
These steps occur at lines~2053--2057 of
Ref.~\cite{Cirac2016MPDO_arXiv}.

For all sufficiently large $L$, the operators
$O_L(A_\gamma)=O_L(M_\gamma)$ are linearly independent. Positivity and the
simple-matrix lemma at lines~1155--1163, applied at line~2058, then give the
multiset identity
\begin{align}
    \{\nu_{\gamma,k}\}_k
        &=
        \left\{\mu_{\alpha,k_1}\mu_{\beta,k_2}
        \chi_{\alpha,\beta,\gamma,k_3}\right\}_{\alpha,\beta,k_1,k_2,k_3}.
\label{eq:cpgsv17_blocked_chi_uniformity_multiplicity_multiset}
\end{align}
Summing the entries in
(\ref{eq:cpgsv17_blocked_chi_uniformity_multiplicity_multiset}) yields
\begin{align}
    \tr(\nu_\gamma)
        &= \sum_{\alpha,\beta}
            m_\alpha m_\beta c^{(1)}_{\alpha,\beta,\gamma},
\label{eq:cpgsv17_blocked_chi_uniformity_two_site_trace_sum}\\
        &= m_\gamma.
\label{eq:cpgsv17_blocked_chi_uniformity_two_site_trace_normalization}
\end{align}
Only after
(\ref{eq:cpgsv17_blocked_chi_uniformity_two_site_trace_normalization}) is
proved do the normalizing maps at lines~2065--2085 of
Ref.~\cite{Cirac2016MPDO_arXiv} take the form
\begin{align}
    \widetilde R_2\!\left(\bigoplus_\gamma X_\gamma\right)
        &= \bigoplus_\gamma
            \frac{\nu_\gamma}{m_\gamma}\otimes X_\gamma,
\label{eq:cpgsv17_blocked_chi_uniformity_two_site_normalizing_map}\\
    \widetilde R_1\!\left(\bigoplus_\gamma X_\gamma\right)
        &= \bigoplus_\gamma
            \frac{\mu_\gamma}{m_\gamma}\otimes X_\gamma.
\label{eq:cpgsv17_blocked_chi_uniformity_one_site_normalizing_map}
\end{align}
The retained-sector composites are
\begin{align}
    T_0 &= \widetilde R_2\widetilde T R_1,
\label{eq:cpgsv17_blocked_chi_uniformity_retained_T_composite}\\
    S_0 &= \widetilde R_1\widetilde S R_2.
\label{eq:cpgsv17_blocked_chi_uniformity_retained_S_composite}
\end{align}
The denominator in
(\ref{eq:cpgsv17_blocked_chi_uniformity_two_site_normalizing_map}) is
normalized by
(\ref{eq:cpgsv17_blocked_chi_uniformity_two_site_trace_normalization}); the
denominator in
(\ref{eq:cpgsv17_blocked_chi_uniformity_one_site_normalizing_map}) is
normalized by the defining identity $m_\gamma=\tr(\mu_\gamma)$. If a
vertical canonical form contains a discarded zero sector, the composites in
(\ref{eq:cpgsv17_blocked_chi_uniformity_retained_T_composite}) and
(\ref{eq:cpgsv17_blocked_chi_uniformity_retained_S_composite}) are trace
preserving only on the retained physical supports. Their completion to the
full physical matrix algebras is documented in
\path{cpgsv17_vertical_isometry_zero_sector.tex}.

The predicate
\leanid{MPOTensor.HasBlockedAdjointFixedPointAlgebraTower} contains neither
the BNT-label coefficients nor the positive
$\chi_{\alpha,\beta,\gamma}$-family or the length-one idempotent trace
identity. It proves
(\ref{eq:cpgsv17_blocked_chi_uniformity_fixed_point_tower}), not the BNT
comparison
(\ref{eq:cpgsv17_blocked_chi_uniformity_source_bnt_comparison}). A
source-faithful construction must first obtain the BNT-label theorem witness,
then use
(\ref{eq:cpgsv17_blocked_chi_uniformity_two_site_trace_normalization}) to
construct the retained-sector composites, and finally complete them on the
discarded physical complements. The criterion
\leanid{MPOTensor.transferRetractData_one_iff_isZCL} concerns the separate
transfer-map idempotence predicate and cannot construct these physical maps.

The scalar normalization is isolated from the still-missing unconditional
channel construction. The structure
\leanid{MPOTensor.BNTMultiplicitySpectrumComparison} records the one-site
traces and the sectorwise multiset equality obtained from the vertical BNT
comparison at lines~2048--2058 of
Ref.~\cite{Cirac2016MPDO_arXiv}. Given this comparison, the BNT algebra clause
proves
\begin{align}
    \tr(\nu_\gamma)
        &= m_\gamma
\label{eq:cpgsv17_blocked_chi_uniformity_formal_trace_normalization}
\end{align}
through
\leanid{MPOTensor.BNTAlgebraClause.twoMultiplicityTrace_eq_traceScalar}.
The idempotent law therefore discharges the scalar obligation at
lines~2058--2064. It neither constructs the comparison from the MPDO tensor
nor defines the maps $T$ and $S$.

The constructor
\leanid{MPOTensor.BNTMultiplicitySpectrumComparison.ofEventuallyEqualPowerSums}
isolates the power-sum step. It proves that positivity and eventual sectorwise
equality of power sums imply the required multiset equality. The proof
combines the geometric-sequence extrapolation used in the BNT Fundamental
Theorem with the unequal-cardinality power-sum identity and formalizes the
argument at lines~2053--2058 of Ref.~\cite{Cirac2016MPDO_arXiv}.

The tensor-level comparison is supplied by the literal two-site spectrum
comparison and exact sector-gauge construction documented in
\path{cpsv16_two_site_sector_unitary_gauge_gap.tex}. The predicate
\leanid{MPOTensor.BNTAlgebraTensorClause} chooses the one-site vertical
canonical decomposition and identifies the abstract operators and trace
scalars with
\begin{align}
    O_L(M_\alpha)
        &\quad\text{and}\quad
    m_\alpha=\sum_q\omega_{\alpha,q},
\label{eq:cpgsv17_blocked_chi_uniformity_tensor_clause_identification}
\end{align}
respectively. It does not include a second vertical canonical decomposition.
Instead, the mixed-prefix argument constructs the relabelling and simultaneous
unitary identifications from the tensor-attached clause. The resulting
positive multiset equality supplies the spectrum comparison without an extra
mathematical hypothesis. The declaration
\leanid{MPOTensor.BNTAlgebraTensorClause.isRFPViaTS} uses this result.

\section{BNT-label coefficient objects}

A separate object represents the same-length coefficient family
$c^{(L)}_{\alpha,\beta,\gamma}$ with fixed BNT labels.\footnote{Formal
declaration: \leanid{MPOTensor.BNTLabelCoefficientFamily}.} If Appendix~C.4
of Ref.~\cite{Cirac2016MPDO_arXiv} has produced a diagonal chi family, the
canonical coefficients are
\begin{align}
    c^{(L)}_{\alpha,\beta,\gamma}
        &= \sum_r\chi_{\alpha,\beta,\gamma,r}^{\,L}.
\label{eq:cpgsv17_blocked_chi_uniformity_canonical_chi_coefficients}
\end{align}
This construction keeps the chi family as the single source of the
coefficients.\footnote{Formal declarations:
\leanid{MPOTensor.BNTLabelCoefficientFamily.ofChi},
\leanid{MPOTensor.BNTLabelCoefficientFamily.ofChi_coeff},
\leanid{MPOTensor.BNTLabelCoefficientFamily.ofChi_coeff_eq_trace_matrix_pow},
and
\leanid{MPOTensor.BNTLabelCoefficientFamily.ofChi_hasPositiveLengthChiTracePowerForm}.}

A same-length BNT product predicate records
\begin{align}
    O_L(M_\alpha)O_L(M_\beta)
        &= \sum_\gamma
            c^{(L)}_{\alpha,\beta,\gamma}O_L(M_\gamma).
\label{eq:cpgsv17_blocked_chi_uniformity_formal_same_length_product}
\end{align}
\footnote{Formal declaration:
\leanid{MPOTensor.BNTLabelOperatorFamily.HasSameLengthProductForm}.}
A predicate on the trace scalars $m_\alpha=\tr(\mu_\alpha)$ records
\begin{align}
    m_\gamma
        &= \sum_{\alpha,\beta}
            c^{(1)}_{\alpha,\beta,\gamma}m_\alpha m_\beta.
\label{eq:cpgsv17_blocked_chi_uniformity_formal_idempotent_condition}
\end{align}
\footnote{Formal declaration:
\leanid{MPOTensor.BNTLabelTraceScalarFamily.HasIdempotentCoefficientForm}.}
A positive BNT-label chi witness consists of a length-independent diagonal
family $\chi_{\alpha,\beta,\gamma}$, positivity of every diagonal entry, and
the positive-length trace-power identity.\footnote{Formal declaration:
\leanid{MPOTensor.PositiveBNTLabelChiTracePowerForm}.} If the coefficients
are canonically determined by the same chi family, positivity alone produces
this witness.\footnote{Formal declaration:
\leanid{MPOTensor.PositiveBNTLabelChiTracePowerForm.ofChi}.}

A blocked-basis comparison predicate states that the chosen blocked
multiplication coefficients are pullbacks of the source-length BNT-label
coefficients. It uses a source label map from the basis of $\mathcal A_n$ and
a target label map from the basis of $\mathcal A_{2n}$. The label assignment
is separate from the coefficient equality, so the Appendix~C.3 construction
target can be stated without assuming that equality.\footnote{Formal
declarations:
\leanid{MPOTensor.BNTBlockedBasisLabelAssignment},
\leanid{MPOTensor.BNTBlockedBasisLabelAssignment.blockedLabel}, and
\leanid{MPOTensor.BNTBlockedBasisCoefficientComparison}.}
This comparison retains the blocked-basis scope restriction: its left-hand
side is expanded in a basis of $\mathcal A_{2n}$, while the source theorem
uses the same-length product
(\ref{eq:cpgsv17_blocked_chi_uniformity_source_same_length_product}). It is
therefore not itself Theorem~4.14(ii) of
Ref.~\cite{Cirac2016MPDO_arXiv}.

The generic declarations do not construct tensor-attached data. That
construction is supplied by
\leanid{MPOTensor.BNTAlgebraTensorClause} and the reflected-target argument
in
\path{TNLean/MPS/MPDO/BNTAlgebraTensorClauseReflectedTarget.lean}. Given the
same-length product form and a positive BNT-label chi witness, one obtains
\begin{align}
    O_L(M_\alpha)O_L(M_\beta)
        &= \sum_\gamma
            \tr\!\left(\chi_{\alpha,\beta,\gamma}^{L}\right)
            O_L(M_\gamma).
\label{eq:cpgsv17_blocked_chi_uniformity_same_length_chi_product}
\end{align}
\footnote{Formal declarations:
\leanid{MPOTensor.BNTLabelOperatorFamily.HasSameLengthProductForm.eq_sum_chi_trace_pow}
and
\leanid{MPOTensor.BNTLabelOperatorFamily.HasSameLengthProductForm.eq_sum_ofChi_trace_pow}.}

The comparison predicate and a positive BNT-label chi witness similarly imply
the blocked coefficient trace-power formula.\footnote{Formal declarations:
\leanid{MPOTensor.BNTBlockedBasisCoefficientComparison.blocked_coeff_eq_trace_pow},
\leanid{MPOTensor.BNTBlockedBasisCoefficientComparison.blocked_coeff_eq_ofChi_trace_pow},
and
\leanid{MPOTensor.BNTBlockedBasisCoefficientComparison.blocked_coeff_eq_pulledBlockedChi_trace_pow}.}
Pulling the uniform BNT-label chi family back along the comparison maps gives a
positive blocked-basis chi witness.\footnote{Formal declarations:
\leanid{MPOTensor.BNTBlockedBasisCoefficientComparison.pulledBlockedChiFamily_toDiagonal_of_pos},
\leanid{MPOTensor.BNTBlockedBasisCoefficientComparison.pulledBlockedChi_tracePowerCoeff_of_pos},
\leanid{MPOTensor.BNTBlockedBasisCoefficientComparison.pulledBlockedChi_dim_of_pos},
\leanid{MPOTensor.BNTBlockedBasisCoefficientComparison.pulledBlockedChi_trace_matrix_pow_of_pos},
and
\leanid{MPOTensor.BNTBlockedBasisCoefficientComparison.toPositiveBlockedStructureChiTracePowerForm}.}
The idempotent condition and the positive chi witness also rewrite the
length-one identity using chi traces.\footnote{Formal declarations:
\leanid{MPOTensor.BNTLabelTraceScalarFamily.HasIdempotentCoefficientForm.eq_sum_chi_trace}
and
\leanid{MPOTensor.BNTLabelTraceScalarFamily.HasIdempotentCoefficientForm.eq_sum_ofChi_trace}.}

The record \leanid{MPOTensor.BNTLabelTheoremData} gathers the fixed
coefficients, same-length operators and product law, trace scalars and
idempotent law, positive BNT-label chi witness, and blocked-basis comparison.
It is not an additional source hypothesis. It packages the data supplied by
the Appendix~C.3--C.4 construction of
Ref.~\cite{Cirac2016MPDO_arXiv}. The canonical constructor
\leanid{MPOTensor.BNTLabelTheoremData.ofChi} covers the special case in which
the chi family determines the coefficients, while retaining the product law,
idempotent law, and blocked-basis comparison as inputs.

The record exposes the source predicates and their equation-level
consequences: the same-length product law, the idempotent scalar law, chi
positivity and trace-power predicates, coefficient trace powers, the
same-length chi-trace product law, the chi-trace idempotent law, blocked-basis
comparison and trace powers, and the positive pulled-back blocked chi
family.\footnote{Formal declarations:
\leanid{MPOTensor.BNTLabelTheoremData.same_length_product_form},
\leanid{MPOTensor.BNTLabelTheoremData.idempotent_coefficient_form},
\leanid{MPOTensor.BNTLabelTheoremData.positive_chi_pos_entries},
\leanid{MPOTensor.BNTLabelTheoremData.positive_chi_trace_power},
\leanid{MPOTensor.BNTLabelTheoremData.labelAssignment},
\leanid{MPOTensor.BNTLabelTheoremData.sourceLabel},
\leanid{MPOTensor.BNTLabelTheoremData.targetLabel},
\leanid{MPOTensor.BNTLabelTheoremData.same_length_product_eq_sum},
\leanid{MPOTensor.BNTLabelTheoremData.idempotent_eq_sum},
\leanid{MPOTensor.BNTLabelTheoremData.coeff_eq_trace_pow},
\leanid{MPOTensor.BNTLabelTheoremData.coeff_eq_ofChi_coeff},
\leanid{MPOTensor.BNTLabelTheoremData.same_length_product_form_ofChi},
\leanid{MPOTensor.BNTLabelTheoremData.idempotent_coefficient_form_ofChi},
\leanid{MPOTensor.BNTLabelTheoremData.same_length_product_eq_sum_ofChi},
\leanid{MPOTensor.BNTLabelTheoremData.idempotent_eq_sum_ofChi},
\leanid{MPOTensor.BNTLabelTheoremData.chi_entry_pos},
\leanid{MPOTensor.BNTLabelTheoremData.blocked_coeff_eq},
\leanid{MPOTensor.BNTLabelTheoremData.blockedComparison_ofChi},
\leanid{MPOTensor.BNTLabelTheoremData.blocked_coeff_eq_ofChi},
\leanid{MPOTensor.BNTLabelTheoremData.same_length_product_eq_sum_chi_trace_pow},
\leanid{MPOTensor.BNTLabelTheoremData.idempotent_eq_sum_chi_trace},
\leanid{MPOTensor.BNTLabelTheoremData.blocked_coeff_eq_trace_pow},
\leanid{MPOTensor.BNTLabelTheoremData.toPositiveBlockedStructureChiTracePowerForm},
\leanid{MPOTensor.BNTLabelTheoremData.positiveBlockedChi},
\leanid{MPOTensor.BNTLabelTheoremData.positiveBlockedChi_toDiagonal_of_pos},
\leanid{MPOTensor.BNTLabelTheoremData.positiveBlockedChi_tracePowerCoeff_of_pos},
\leanid{MPOTensor.BNTLabelTheoremData.positiveBlockedChi_dim_of_pos},
\leanid{MPOTensor.BNTLabelTheoremData.positiveBlockedChi_trace_matrix_pow_of_pos},
\leanid{MPOTensor.BNTLabelTheoremData.positiveBlockedChi_tracePower},
\leanid{MPOTensor.BNTLabelTheoremData.positiveBlockedChi_posEntries},
\leanid{MPOTensor.BNTLabelTheoremData.positiveBlockedChi_entry_pos}, and
\leanid{MPOTensor.BNTLabelTheoremData.blocked_coeff_eq_positiveBlockedChi_trace_pow}.}

The existential form closest to the source is
\leanid{MPOTensor.BNTLabelTheoremWitness}. It records the BNT-label type,
same-length operator spaces, algebraic structure, and theorem data. The
predicate \leanid{MPOTensor.HasBNTLabelTheoremWitness} asserts that such a
witness exists. It can be unpacked into fixed coefficients, operators, trace
scalars, a positive chi family, blocked comparison maps, the source product
and idempotent laws, their chi-trace forms, and the pulled-back blocked chi
family.\footnote{Formal declarations:
\leanid{MPOTensor.BNTLabelTheoremWitness.toTheoremData},
\leanid{MPOTensor.BNTLabelTheoremWitness.coeffs},
\leanid{MPOTensor.BNTLabelTheoremWitness.operators},
\leanid{MPOTensor.BNTLabelTheoremWitness.traceScalars},
\leanid{MPOTensor.BNTLabelTheoremWitness.positiveChi},
\leanid{MPOTensor.BNTLabelTheoremWitness.blockedComparison},
\leanid{MPOTensor.BNTLabelTheoremWitness.same_length_product_form},
\leanid{MPOTensor.BNTLabelTheoremWitness.idempotent_coefficient_form},
\leanid{MPOTensor.BNTLabelTheoremWitness.positive_chi_pos_entries},
\leanid{MPOTensor.BNTLabelTheoremWitness.positive_chi_trace_power},
\leanid{MPOTensor.BNTLabelTheoremWitness.labelAssignment},
\leanid{MPOTensor.BNTLabelTheoremWitness.sourceLabel},
\leanid{MPOTensor.BNTLabelTheoremWitness.targetLabel}, and
\leanid{MPOTensor.BNTLabelTheoremWitness.positiveBlockedChi}.}

The constructor
\leanid{MPOTensor.BNTLabelTheoremWitness.ofChi} gives the analogous
existential object from a positive chi family. It still takes the product,
idempotent, and blocked-comparison statements as inputs; it does not derive
them from an arbitrary MPDO tensor. The witness then yields the coefficient
trace-power formula, product and idempotent laws, positivity, blocked
comparison, chi-trace rewrites, and positive pulled-back blocked
data.\footnote{Formal declarations:
\leanid{MPOTensor.BNTLabelTheoremWitness.same_length_product_form},
\leanid{MPOTensor.BNTLabelTheoremWitness.idempotent_coefficient_form},
\leanid{MPOTensor.BNTLabelTheoremWitness.positive_chi_pos_entries},
\leanid{MPOTensor.BNTLabelTheoremWitness.positive_chi_trace_power},
\leanid{MPOTensor.BNTLabelTheoremWitness.sourceLabel},
\leanid{MPOTensor.BNTLabelTheoremWitness.targetLabel},
\leanid{MPOTensor.BNTLabelTheoremWitness.coeff_eq_trace_pow},
\leanid{MPOTensor.BNTLabelTheoremWitness.coeff_eq_ofChi_coeff},
\leanid{MPOTensor.BNTLabelTheoremWitness.same_length_product_form_ofChi},
\leanid{MPOTensor.BNTLabelTheoremWitness.idempotent_coefficient_form_ofChi},
\leanid{MPOTensor.BNTLabelTheoremWitness.same_length_product_eq_sum_ofChi},
\leanid{MPOTensor.BNTLabelTheoremWitness.idempotent_eq_sum_ofChi},
\leanid{MPOTensor.BNTLabelTheoremWitness.same_length_product_eq_sum},
\leanid{MPOTensor.BNTLabelTheoremWitness.idempotent_eq_sum},
\leanid{MPOTensor.BNTLabelTheoremWitness.chi_entry_pos},
\leanid{MPOTensor.BNTLabelTheoremWitness.blocked_coeff_eq},
\leanid{MPOTensor.BNTLabelTheoremWitness.blockedComparison_ofChi},
\leanid{MPOTensor.BNTLabelTheoremWitness.blocked_coeff_eq_ofChi},
\leanid{MPOTensor.BNTLabelTheoremWitness.same_length_product_eq_sum_chi_trace_pow},
\leanid{MPOTensor.BNTLabelTheoremWitness.idempotent_eq_sum_chi_trace},
\leanid{MPOTensor.BNTLabelTheoremWitness.blocked_coeff_eq_trace_pow},
\leanid{MPOTensor.BNTLabelTheoremWitness.toPositiveBlockedStructureChiTracePowerForm},
\leanid{MPOTensor.BNTLabelTheoremWitness.positiveBlockedChi_toDiagonal_of_pos},
\leanid{MPOTensor.BNTLabelTheoremWitness.positiveBlockedChi_tracePowerCoeff_of_pos},
\leanid{MPOTensor.BNTLabelTheoremWitness.positiveBlockedChi_dim_of_pos},
\leanid{MPOTensor.BNTLabelTheoremWitness.positiveBlockedChi_trace_matrix_pow_of_pos},
\leanid{MPOTensor.HasBNTLabelTheoremWitness.positive_blocked_chi_witness},
\leanid{MPOTensor.HasBNTLabelTheoremWitness.exists_blocked_chi_trace_power_form},
\leanid{MPOTensor.HasBNTLabelTheoremWitness.exists_blocked_coeff_eq_trace_pow},
\leanid{MPOTensor.BNTLabelTheoremWitness.positiveBlockedChi_tracePower},
\leanid{MPOTensor.BNTLabelTheoremWitness.positiveBlockedChi_posEntries},
\leanid{MPOTensor.BNTLabelTheoremWitness.positiveBlockedChi_entry_pos}, and
\leanid{MPOTensor.BNTLabelTheoremWitness.blocked_coeff_eq_positiveBlockedChi_trace_pow}.}

Neither an arbitrary MPDO tensor nor the support-algebra predicate alone can
produce this witness. Such a construction would imply the false converse
demonstrated by the phase-flip channel. In Theorem~4.14(ii) of
Ref.~\cite{Cirac2016MPDO_arXiv}, the BNT-label coefficients, their positive
trace-power representation, and the length-one idempotent law are hypotheses
of the converse rather than consequences of stabilized adjoint fixed-point
algebras.

\section{Length-independent zipper specialization}

Consider a fusion-isometry family satisfying Theorem~4.14(iii) of
Ref.~\cite{Cirac2016MPDO_arXiv}. Assume that its positive trace-power
coefficients are independent of the positive chain length. The observation at
line~1010 of that source gives
\begin{align}
    \chi_{\alpha,\beta,\gamma}
        &= 1_{r_{\alpha,\beta,\gamma}}.
\label{eq:cpgsv17_blocked_chi_uniformity_length_independent_chi}
\end{align}
Substitution into the weighted fusion identity yields both zipper equations
and the completeness identity
\begin{align}
    (M_\alpha M_\beta)^{ij}
        &=
        U_{\alpha,\beta}^\dagger
        \left(\bigoplus_\gamma
            1_{r_{\alpha,\beta,\gamma}}\otimes M_\gamma^{ij}\right)
        U_{\alpha,\beta}.
\label{eq:cpgsv17_blocked_chi_uniformity_unweighted_fusion_identity}
\end{align}
These conclusions are formalized by
\leanid{MPOTensor.BNTFusionIsometryFamily.fusionIsometry_mul_mulTensor_of_lengthIndependent},
\leanid{MPOTensor.BNTFusionIsometryFamily.mulTensor_mul_fusionIsometry_conjTranspose_of_lengthIndependent},
and
\leanid{MPOTensor.BNTFusionIsometryFamily.mulTensor_eq_conjTranspose_mul_unweightedDirectSum_mul_of_lengthIndependent}.
They are the identity-weight specialization of
\texttt{zippercondition2} and \texttt{inversegaugeone} in the ``Fusion
tensors'' section of the arXiv:1511.08090 source file
\texttt{AnyonsPEPS.tex}.

This statement is scope restricted: it assumes the length-independent
coefficient system discussed at lines~995--1010 of
Ref.~\cite{Cirac2016MPDO_arXiv}; it does not prove that every renormalization
fixed point has such coefficients. In this specialization, canonical bond
reassociation and the two iterated fusion isometries compare the complete
left- and right-associated direct sums. The formal letterwise identity says
that this comparison intertwines the two unweighted direct sums. The
comparison matrix is oriented from the right-associated sum to the
left-associated sum, opposite to the printed $F$-matrix convention in
\texttt{Fmove}; it is not identified with that printed matrix.

Let $C$ be the full comparison, and let $U^{\mathrm L}$ and
$U^{\mathrm R}$ be the left- and right-associated iterated fusion
isometries. Then
\begin{align}
    CC^\dagger
        &= U^{\mathrm L}(U^{\mathrm L})^\dagger,
\label{eq:cpgsv17_blocked_chi_uniformity_left_range_projection}\\
    C^\dagger C
        &= U^{\mathrm R}(U^{\mathrm R})^\dagger.
\label{eq:cpgsv17_blocked_chi_uniformity_right_range_projection}
\end{align}
These identities are formalized by
\leanid{MPOTensor.BNTFusionIsometryFamily.tripleFusionComparison_mul_conjTranspose}
and
\leanid{MPOTensor.BNTFusionIsometryFamily.conjTranspose_mul_tripleFusionComparison}.
The fusion assumptions give
\begin{align}
    (U^{\mathrm L})^\dagger U^{\mathrm L}
        &= 1,
\label{eq:cpgsv17_blocked_chi_uniformity_left_isometry}\\
    (U^{\mathrm R})^\dagger U^{\mathrm R}
        &= 1.
\label{eq:cpgsv17_blocked_chi_uniformity_right_isometry}
\end{align}
Without further information, $C$ is only a partial isometry between the two
ranges. The invertible total fusion matrix at lines~187--191 of the
arXiv:1511.08090 source supplies the corresponding completeness in the source
argument.

The simultaneous inverse relation at line~269 of that source is
\begin{align}
    B_d^+B_{d'}
        &= \delta_{d,d'}1.
\label{eq:cpgsv17_blocked_chi_uniformity_simultaneous_inverse_relation}
\end{align}
It separates all final-label blocks at once. The predicate
\leanid{MPOTensor.BNTFusionIsometryFamily.HasFinalLabelSelectorWords}
records its exact finite-word consequence as a hypothesis. At a positive word
length, the selector polynomial supplies the missing completeness: every
nonempty word in the unweighted direct sum lies in the range of the
corresponding fusion isometry, and summing the selector polynomials over the
final label gives the identity. Consequently,
\begin{align}
    U^{\mathrm L}(U^{\mathrm L})^\dagger
        &= 1,
\label{eq:cpgsv17_blocked_chi_uniformity_left_coisometry}\\
    U^{\mathrm R}(U^{\mathrm R})^\dagger
        &= 1,
\label{eq:cpgsv17_blocked_chi_uniformity_right_coisometry}\\
    CC^\dagger
        &= 1,
\label{eq:cpgsv17_blocked_chi_uniformity_comparison_left_unitarity}\\
    C^\dagger C
        &= 1.
\label{eq:cpgsv17_blocked_chi_uniformity_comparison_right_unitarity}
\end{align}
These statements are formalized by
\leanid{MPOTensor.BNTFusionIsometryFamily.fusionIsometry_mul_conjTranspose_eq_one_of_lengthIndependent_of_selectorWords},
\leanid{MPOTensor.BNTFusionIsometryFamily.leftFusionIsometry_mul_conjTranspose_eq_one_of_lengthIndependent_of_selectorWords},
\leanid{MPOTensor.BNTFusionIsometryFamily.rightFusionIsometry_mul_conjTranspose_eq_one_of_lengthIndependent_of_selectorWords},
\leanid{MPOTensor.BNTFusionIsometryFamily.tripleFusionComparison_mul_conjTranspose_eq_one_of_lengthIndependent_of_selectorWords},
and
\leanid{MPOTensor.BNTFusionIsometryFamily.conjTranspose_mul_tripleFusionComparison_eq_one_of_lengthIndependent_of_selectorWords}.
The positive-length restriction is necessary because the telescoping argument
uses at least one conjugated letter.

Under the same selector hypothesis,
\leanid{MPOTensor.BNTFusionIsometryFamily.tripleFusionComparison_finalSector_submatrix_eq_zero}
proves that the comparison has a zero corner between distinct final labels.
The proof takes a rectangular bond-space corner of the letterwise
intertwining identity, extends it to words, and applies the selector
polynomial.

Assume additionally that the selected final tensor is injective at the
current blocking. Then
\leanid{MPOTensor.BNTFusionIsometryFamily.exists_tripleFusionComparison_finalSector_eq_kronecker_one_of_separation}
combines sector separation with the amplified-commutant argument. For every
selected final label $\varepsilon$, it produces a rectangular multiplicity
matrix $F_\varepsilon$ such that
\begin{align}
    C^{\varepsilon,\varepsilon}_{\alpha,\beta,\gamma}
        &= F_\varepsilon\otimes 1_{D_\varepsilon},
\label{eq:cpgsv17_blocked_chi_uniformity_final_sector_factorization}\\
    C^{\varepsilon,\varepsilon'}_{\alpha,\beta,\gamma}
        &= 0,
&
    \varepsilon'&\ne\varepsilon.
\label{eq:cpgsv17_blocked_chi_uniformity_offdiagonal_final_sector}
\end{align}
This is the bond-factor conclusion at line~277 of the arXiv:1511.08090
source, with the opposite comparison orientation.

For an abstract fusion-isometry family, this remains a scope restriction. If
the labelled tensors are linked to the source BNT witness, simultaneous
final-label selectors are derived rather than assumed. One-letter injectivity
at the chosen blocking remains separate, and the factorization theorem alone
permits $F_\varepsilon$ to be rectangular.

Length independence and eventual linear independence also imply the
associativity identity
\begin{align}
    \sum_\delta
        r_{\alpha,\beta}^{\delta}
        r_{\delta,\gamma}^{\varepsilon}
        &=
        \sum_\delta
        r_{\beta,\gamma}^{\delta}
        r_{\alpha,\delta}^{\varepsilon}.
\label{eq:cpgsv17_blocked_chi_uniformity_multiplicity_associativity}
\end{align}
This is formalized by
\leanid{MPOTensor.BNTFusionIsometryFamily.card_leftFinalMultiplicity_eq_card_rightFinalMultiplicity_of_lengthIndependent}.
It is the equality of fixed-final multiplicity dimensions at lines~238--241
of the arXiv:1511.08090 source. It proves that $F_\varepsilon$ is square when
the factorization and eventual-independence hypotheses both hold, but it does
not by itself prove invertibility.

The selector hypothesis removes the range-completeness obstruction from
lines~181--191 of that source. The two-sided unitarity of $C$, together with
the vanishing off-diagonal final-label corners, gives two-sided unitarity on
each fixed-final bond-space corner. The corresponding formal statements are
\leanid{MPOTensor.BNTFusionIsometryFamily.tripleFusionComparison_finalSector_mul_conjTranspose_eq_one}
and
\leanid{MPOTensor.BNTFusionIsometryFamily.conjTranspose_mul_tripleFusionComparison_finalSector_eq_one}.

Combining these corner identities with
(\ref{eq:cpgsv17_blocked_chi_uniformity_final_sector_factorization}) gives
the two unitary identities for $F_\varepsilon$ after cancelling the identity
factor. The declaration
\leanid{MPOTensor.BNTFusionIsometryFamily.exists_tripleFusionComparison_finalSector_eq_kronecker_one_and_unitary}
formalizes this conclusion under $D_\varepsilon>0$. This qualification is
necessary because injectivity is vacuous when $D_\varepsilon=0$. The
fixed-final invertibility obstruction is therefore removed under the
simultaneous-selector, present-blocking injectivity, and
positive-bond-dimension hypotheses.

For a general coefficient-form BNT,
\leanid{MPSTensor.IsCPSVBasisOfNormalTensors.exists_blocked_wordTupleSpanTop_one}
derives positive bond dimensions, a common positive blocking length, and the
simultaneous one-letter span of the blocked representatives. The span contains
every tuple that is the identity on one labelled block and zero on all others.
Hence one-letter injectivity and selector coefficients can be extracted after
the same common blocking without adding either property to the BNT definition.

The pre-blocking common word span is exposed by
\leanid{MPSTensor.IsCPSVBasisOfNormalTensors.exists_positive_wordTupleSpanTop}.
Its selector coefficients connect directly to the fusion-family formulation.
Consequently,
\leanid{MPOTensor.BNTFusionIsometryFamily.fusionIsometry_mul_conjTranspose_eq_one_of_bnt_of_lengthIndependent}
proves
\begin{align}
    U_{\alpha,\beta}U_{\alpha,\beta}^{\dagger}
        &= I
\label{eq:cpgsv17_blocked_chi_uniformity_bnt_derived_coisometry}
\end{align}
from the BNT witness and length-independent fusion data, without a selector
or coisometry hypothesis. The BNT argument supplies a positive selector
length, so the empty-chain case is unused. This is a short consequence of the
source ingredients, not a proof printed verbatim in either source.

For the positive-diagonal fusion family, assume explicitly that the labelled
doubled-index tensors are identified with a BNT witness. The declaration
\leanid{MPOTensor.BNTFusionIsometryFamily.exists_positive_block_with_injective_fusion}
transports the construction through the same positive block length. It
derives the common one-letter span and injectivity of every blocked labelled
tensor and proves the blocked conjugation identity, both zipper orientations,
and literal reconstruction. Its multiplicity is the dimension of the
specialized chi space. The blocking length is positive; the empty-chain
conjugation formula is not used to prove coisometry, which instead follows
from positive-length BNT selectors.

The declaration
\leanid{MPSTensor.WordTupleSpanTop.exists_mpo_block_left_inverse} extracts the
simultaneous block-letter inverse from the same one-letter product-algebra
span. After reindexing the product bond space as a pair,
$U_{\alpha,\beta}^\dagger$ is the synthesis map and $U_{\alpha,\beta}$ is the
analysis map. These maps and the blocked identities are assembled by
\leanid{MPOTensor.BNTFusionIsometryFamily.exists_blocked_completeZipperFusionFamily}.
Thus the existing printed $F$-matrix and pentagon theorem apply to the
length-independent source BNT fusion family without an additional
injectivity, selector, inverse, or coisometry hypothesis.

The RFP tensor supplies both the labelled fusion-isometry family and its
identification with the BNT of the vertical canonical form through
\leanid{MPOTensor.HasBNTFusionTensorClause.of_isRFPViaTS}. This closes the
corrected active-support form of the Appendix~C.3--C.4 implication in
Ref.~\cite{Cirac2016MPDO_arXiv}. It is logically independent of the complete
zipper construction that follows.

The source-faithful complete zipper structure is formalized separately by
\leanid{MPOTensor.CompleteZipperFusionFamily}. Its data consist of integer
fusion multiplicities, independent synthesis and analysis tensors, their
biorthogonality, the product-letter reconstruction identity
\texttt{inversegaugeone}, both zipper equations, injectivity of the labelled
blocks, and the simultaneous one-letter inverse from lines~269--277 of
\texttt{AnyonsPEPS.tex}. The product-MPO projector at lines~181--191 is not
identified with the identity on the ambient product bond space.

Applying \texttt{inversegaugeone} to the two associative decompositions at
lines~238--247 reconstructs each triple-product letter through either fusion
tree. Contracting these equal reconstructions with the simultaneous inverse
gives equality of the two triple-product support projectors. These assumptions
construct
\leanid{MPOTensor.CompleteZipperFusionFamily.printedFMatrix}, whose rows index
right-tree multiplicities and whose columns index left-tree multiplicities.
The exact \texttt{Fmove} identity and both inverse identities are formalized
by
\leanid{MPOTensor.CompleteZipperFusionFamily.rightTripleSynthesis_mul_printedFMatrix},
\leanid{MPOTensor.CompleteZipperFusionFamily.printedFMatrix_mul_inversePrintedFMatrix},
and
\leanid{MPOTensor.CompleteZipperFusionFamily.inversePrintedFMatrix_mul_printedFMatrix}.
No length-independent coefficient system, selector words, or additional
present-blocking injectivity hypothesis is imposed on this source-labelled
result.

The complete zipper structure does not close the remaining gap for the
positive-diagonal family: constructing complete zipper data from the
assumptions of Ref.~\cite{Cirac2016MPDO_arXiv} remains open. Within the
complete zipper family, however, the two-sided printed $F$-move and the
fourfold entrywise pentagon identity from lines~279--283 of
\texttt{AnyonsPEPS.tex} are formalized.

The five categorical fourfold multiplicity spaces, their fusion maps, and the
five lifted printed-$F$ edges are represented by
\leanid{MPOTensor.CompleteZipperFusionFamily.FourfoldLeftAssocMultiplicity}
through
\leanid{MPOTensor.CompleteZipperFusionFamily.twoEdgePrintedFMatrix}. They use
the integer fusion multiplicities $N_{ab}^{c}$ of the complete zipper family
and are not identified with the positive-diagonal weighted coordinates of the
restricted family. The five elementary moves identify the three-edge and
two-edge composites with the same expansion in the right-associated fourfold
fusion map. Its analysis map cancels this expansion, proving
\leanid{MPOTensor.CompleteZipperFusionFamily.threeEdgePrintedFMatrix_eq_twoEdgePrintedFMatrix}.
The inverse edge matrices then give
\leanid{MPOTensor.CompleteZipperFusionFamily.threeEdgeInversePrintedFMatrix_eq_twoEdgeInversePrintedFMatrix}
and the literal indexed equation~(33) of Ref.~\cite{Cirac2016MPDO_arXiv},
\leanid{MPOTensor.CompleteZipperFusionFamily.inversePrintedFMatrix_pentagon}.

\paragraph{Local fixes in the fusion-tensor display.}

The prose at line~161 of \texttt{AnyonsPEPS.tex} points the arrow for
$X^c_{ab,\mu}$ from the product bond space to the $c$ bond space. The
identities at lines~163--200 and the printed $F$-move at lines~248--251
require the reverse typing:
\begin{align}
    X^c_{ab,\mu}
        &\colon \mathbb C^{D_c}
        \longrightarrow
        \mathbb C^{D_a}\otimes\mathbb C^{D_b}.
\label{eq:cpgsv17_blocked_chi_uniformity_fusion_tensor_orientation}
\end{align}
With this orientation, $X^{c+}_{ab,\mu}$ is the left inverse. The same source
display writes the Kronecker factor in
$X^{d+}_{ab,\nu}X^c_{ab,\mu}$ as $\delta_{de}$, although $e$ does not occur
elsewhere in the relation. The typed factor is
\begin{align}
    \delta_{dc}
        &= \delta_{cd}.
\label{eq:cpgsv17_blocked_chi_uniformity_correct_kronecker_delta}
\end{align}
The formalization uses these equation-compatible conventions.

\paragraph{Local fix in the pentagon index orientation.}

The source identity \texttt{Fmove} at lines~248--251 places right-tree indices
above and left-tree indices below:
\begin{align}
    \left(F^{abc}_{d}\right)^{f\lambda\sigma}_{e\mu\nu}.
\label{eq:cpgsv17_blocked_chi_uniformity_printed_F_index_orientation}
\end{align}
Each factor in \texttt{pentagoneq} at lines~280--283 reverses this placement.
With the multiplicity spaces fixed by the fusion trees at line~279, those
entries belong to the inverse $F$-matrix. The forward edge matrices in the
formalization follow \texttt{Fmove}; equivalently, the literal upper and lower
indices in \texttt{pentagoneq} are interpreted through
\leanid{MPOTensor.CompleteZipperFusionFamily.inversePrintedFMatrix}. The
formalization states both the forward matrix identity and the literally
indexed inverse identity.

\section{Completion status for Theorem 4.14}

The tensor-attached BNT route is complete. The mixed one-site/two-site prefix
argument supplies the unitary relabelling and the positive multiset identity
for $\nu_{\gamma,k}$. The retained-sector composites
(\ref{eq:cpgsv17_blocked_chi_uniformity_retained_T_composite}) and
(\ref{eq:cpgsv17_blocked_chi_uniformity_retained_S_composite}), together with
the zero-sector measure-and-prepare completion, give the two trace-preserving
completely positive maps required by Definition~4.1 of
Ref.~\cite{Cirac2016MPDO_arXiv}. The declaration
\leanid{MPOTensor.BNTAlgebraTensorClause.isRFPViaTS} formalizes this
implication.

Under the literal CPSV canonical-form and MPDO hypotheses, the declaration
\leanid{MPOTensor.isRFPViaTS_iff_hasBNTAlgebraTensorClause} combines the two
directions. The equivalence
\leanid{MPOTensor.isRFPViaTS_iff_hasBNTFusionTensorClause} uses the documented
active-support correction of Theorem~4.14(iii), rather than its unrestricted
printed form. These results neither identify the positive-diagonal family with
a complete zipper fusion family nor provide the stronger categorical
$F$-matrix data discussed in the preceding section.

The one-site criterion for \leanid{MPOTensor.TransferRetractData} remains
separate. It is equivalent to doubled-index transfer-map idempotence, not to
Definition~4.1 of Ref.~\cite{Cirac2016MPDO_arXiv}.

\bibliographystyle{alpha}
\bibliography{references}

\end{document}
